% Generated draft: requires mathematical and source review.
\documentclass[11pt]{article}
\usepackage[margin=25mm]{geometry}
\usepackage{fontspec}
\setmainfont{Latin Modern Roman}
\usepackage{amsmath,amssymb,mathtools,unicode-math}
\setmathfont{Latin Modern Math}
\usepackage{enumitem,needspace,xurl,hyperref}
\setlength{\parindent}{0pt}
\setlength{\parskip}{5pt}
\setlength{\emergencystretch}{3em}
\newcommand{\N}{\mathbb{N}}
\newcommand{\Z}{\mathbb{Z}}
\newcommand{\Q}{\mathbb{Q}}
\newcommand{\R}{\mathbb{R}}
\newcommand{\C}{\mathbb{C}}
\newcommand{\F}{\mathbb{F}}
\newcommand{\A}{\mathbb{A}}
\newcommand{\PP}{\mathbb{P}}
\newcommand{\E}{\mathbb{E}}
\newcommand{\eps}{\varepsilon}
\newcommand{\dd}{\,\mathrm{d}}
\begin{document}
\section*{UnsolvedMath ID: 700041}
\textbf{Source problem code:} catalogue entry
\subsection{Definitions and mathematical statement}
Let \(G\) be a group and \(H\leq G\). The normalizer of \(H\) in \(G\) is
\[
N_G(H)=\{g\in G:gHg^{-1}=H\}.
\]
The supplied source does not define ``relatively convex'' or specify an ordering convention. Consequently, the following formulation is conditional on an explicit interpretation.

Under the interpretation using ordered cosets, call \(H\) relatively convex in \(G\) if the set of left cosets \(G/H\) admits a strict total order \(\prec\) invariant under the natural left action of \(G\):
\[
\forall a,x,y\in G,\qquad xH\prec yH\Longrightarrow axH\prec ayH.
\]
The problem is to determine conditions on \(G\) and \(H\) ensuring that, if \(G/H\) admits such an order, then \(G/N_G(H)\) also admits a strict total order invariant under the natural left action of \(G\). The two coset orders need not arise from any specified common order.

Another possible interpretation requires an ambient group order. A strict total order \(<\) on \(G\) is left invariant if
\[
\forall a,x,y\in G,\qquad x<y\Longrightarrow ax<ay.
\]
A subgroup \(K\leq G\) is convex for this order if
\[
\forall k_1,k_2\in K\ \forall g\in G,\qquad k_1<g<k_2\Longrightarrow g\in K.
\]
If ``relatively convex'' means convex for some left-invariant strict total order on \(G\), the question instead asks for conditions ensuring that the existence of such an order making \(H\) convex implies the existence of such an order making \(N_G(H)\) convex. The supplied record does not settle whether this convention, the ordered-coset convention, or a convention requiring a two-sided invariant order is intended. These qualifications require source clarification before a unique formulation can be fixed.
\subsection{Short English statement}
Find conditions under which the normalizer \(N_G(H)\) of a relatively convex subgroup \(H\) is itself relatively convex in \(G\). The source does not specify the meaning of relative convexity or the required ordering convention, so the precise interpretation remains unresolved.
\subsection{Sources}
\begin{itemize}
\item Local UnsolvedMath registry, supplied complete source record: permanent ID \(700041\), exact source problem code \texttt{catalogue entry}, title ``Kourovka Notebook Problem 1.46.''
\item Attributed in the supplied record to A. I. Kokorin, proposed in 1965, original problem number 1.46 in \emph{The Kourovka Notebook (Unsolved Problems in Group Theory)}, first issue (1965). The record cites the twenty-first issue edition (2026), \texttt{arXiv:1401.0300v48}, and identifies E. I. Khukhro and V. D. Mazurov as editors of the collection.
\item The registry labels the problem \texttt{open}. This is supplied source metadata, not an independently verified assessment. The record supplies no literature-verification date, page reference, or additional literature sources; the cited edition has not been independently consulted for this draft.
\end{itemize}
\end{document}
